%% The manuscript body, shared by every binding of this paper.
%% The class, the front matter and the bibliography style belong to the
%% binding; everything a reader reads belongs here, once.
\section{Introduction}

Computational research is routinely sorted into reproducible and not
reproducible, and artifact-evaluation programmes reinforce the binary by awarding
badges that certify code and data were made available~\cite{acmbadging2020}. The
binary is convenient and it is also the wrong shape. Reproduction is a sequence
of steps that can each succeed or fail for unrelated reasons, and a reader who
gets stuck is almost never told which step stopped them. Large-scale attempts to
build the code accompanying published papers have found the sticking points to be
mundane and specific, and to sit at very different depths for different
papers~\cite{collberg2016repeatability}. The literature on the reproducibility of
computational work has argued for a decade that the useful unit is a gradation
rather than a verdict~\cite{peng2011reproducible,baker2016reproducibility,
donoho2010invitation,ince2012open,stodden2016science,goodman2016mean,
plesser2018repro}. That argument now sits inside a wider accounting of why
published findings fail to travel: low power, undisclosed flexibility, and
selective reporting~\cite{ioannidis2005why,ioannidis2014true,button2013power,
simmons2011false,head2015phacking,openscience2015repro,begley2012raise,
nature2016reality,munafo2017manifesto}.

This study takes that argument literally and treats the stopping point as the
object to be recorded. The instrument is a ladder of \NRungs{} rungs, described
in the methods, on which each rung is an assertion checkable on a given machine
without contacting anyone: a citable record exists; the source is obtainable at a
fixed revision; the runtime that source requires is present; the published table
recomputes. Because a break makes everything above it unreachable, the reportable
quantity is the index of the first break. An audit result is therefore a small
integer plus the observation that justifies it, and it is refutable line by line
rather than in aggregate.

We apply the ladder to one published comparison in acoustic sensor calibration
and report what a read-only probe observed. The result is that the break is
neither at the top nor at the bottom. The source is present at exactly the
revision recorded for it, with every one of its \EntryPoints{} entry points on
disk, so the availability that a badge would certify is genuinely there. What is
missing is the commercial numerical environment the source is written for. That
single absence is sufficient, and everything above it follows.

Corpus availability turns out not to fit on the ladder at all, and we report it
as a separate axis for a reason worth stating plainly. The evaluation corpus for
this line of work is publicly archived, openly downloadable and about
\CorpusGB{} gigabytes, and it is not on the machine. Calling that a satisfied
requirement misdescribes the situation, and calling it a broken one blames the
wrong party. A resolvable link is not a local asset, and an audit that collapses
the two will systematically overstate how close a reader is to a reproduction.

The audit leaves a large part of the project fully reproducible, and that part
supplies a second result which is not about reproducibility at all. A local
simulation compares an offset-agnostic estimator against an offset-aware one
across \NConditions{} conditions. At every condition where the unknown offsets it
models are actually present, the offset-aware estimator wins, as expected. At the
single condition where the offsets are zero, it loses by a factor of about
\ZeroRatio{} in median error and loses on every one of \NTrials{} paired trials.
That condition is the one a method paper is least likely to include, and its
omission would leave a headline that is true at \NonZeroConditions{} of the
\NConditions{} conditions and badly false at the remaining one.

A third record in the same project is the one that changed what we think the
ladder is for. It compares two locally written estimators over \CxPlanned{}
paired trials, and it stores both the cell summary a paper would print and the
individual trials behind it. Read through the mean, it names one estimator the
better at \CxDisagreeing{} of \CxConditions{} conditions; read through the
median, or through a paired test on the same trials, it names the other at all
\CxConditions{}. At one of those conditions every planned trial completed in
both arms, so no handling of missing runs is involved and the disagreement is
the summary and nothing else. Setting aside \CxFlipDepth{} trial in
\CxTrialsPerCondition{} is enough to make the mean agree.

That is a reader's problem, not an authors' one, and it is invisible to the
ladder. Rung \NRungs{} asks whether the published table recomputes; a reader who
reaches it has reproduced a number, and the number does not determine the claim.
Recomputing a table of means would faithfully reproduce a conclusion that the
trials underneath it do not support. So we record a second off-ladder axis
beside corpus availability --- whether the per-trial record is disclosed at all
--- and we report what disclosure would have weighed here: fewer bytes than the
summary it would have accompanied.

The fourth result is the one that shows disclosure is not sufficient either, and
it required an experiment rather than a re-reading. The offset study attaches
its reversal to a single condition because its grid steps from zero to
\CoarseGapHigh{} milliseconds in one move and has nothing in between. We reran
that comparison on a grid refined inside the skipped interval and crossed with
\FineRooms{} independently simulated arrays, \FinePairs{} paired trials in all,
with the estimator code unchanged. The reversal is not a point. The
offset-agnostic estimator holds the lower median error out to between
\ReversalLow{} and \ReversalHigh{} milliseconds depending on the array, and at
the original sweep's lowest non-zero condition it is still the better of the two
in \AnchorAgnosticRooms{} of \FineRooms{} arrays. Where the ordering changes sign
is therefore a joint property of the estimators and the grid, and no summary of
the conditions a study ran can recover the conditions it did not. Disclosing
trials cannot repair that, because the trials do not exist; what can is one more
condition, and the record examined here contains enough to say where to put it.

\section{Related work}

\subsection{Reproducibility as a verdict, and the attempts to grade it}

The position that computational claims should travel with the code and data that
produced them is old enough to have infrastructure built around
it~\cite{claerbout1992electronic,fomel2009reproducible,peng2011reproducible,
donoho2010invitation,ince2012open,vandewalle2009repro,stodden2014practices,
sandve2013ten,wilson2014best,wilson2017good,gundersen2021fundamental}. The
infrastructure that arrived was largely binary. Artifact-evaluation programmes
award badges for availability, functionality and
reusability~\cite{acmbadging2020}, machine-learning venues have run reproduction
challenges with the same threshold shape~\cite{pineau2021improving}, and each
badge is a threshold: an artifact either clears it or it does not, and the
reader who fails to reproduce something learns nothing about where they stopped.
Surveys of researchers report a widespread sense that published results often
cannot be reconstructed~\cite{baker2016reproducibility,nature2016reality,
hutson2018ai,gundersen2018state}, without saying at which step the
reconstruction fails, because the survey instrument cannot ask that. Attempts to
re-run published computational science as a community
practice~\cite{rougier2017rescience,landis2012transparent,nosek2015culture,
benureau2018rerun,haibekains2020transparency,heil2021reproducibility}
likewise grade a finished artifact, not a stopping point.

The studies that did try to run the code found the answer to be specific and
mundane. Collberg and Proebsting attempted to build the artifacts of several
hundred systems papers and classified the outcomes by how far they
got~\cite{collberg2016repeatability}; the classification is close in spirit to
what we do here, and the difference is that we make the ordering a stated
property of the instrument rather than a convenient way of presenting results,
and we require each level to be an assertion refutable on a named machine on a
named date. Proposition~\ref{prop:sufficient} is what that discipline buys: the
observation collapses to an integer without loss, which a set of badges does
not.

Work on capturing the execution environment --- containers, declarative
environment specifications, workflow systems --- attacks the same problem from
the other end, by trying to make rung three hold by
construction~\cite{boettiger2015docker,kurtzer2017singularity,
stodden2014practices,sandve2013ten}.
It is complementary to an audit and does not replace one: a container image is
itself an artifact that can be absent, stale, or built for an architecture the
reader does not have, and the question of where a given reader stops remains
empirical. The break we record here, a commercial numerical runtime the source
is written for, is also the case those tools address least well, since the
licence is not something an image can carry.

\subsection{Availability, obtainability, and the distance between them}

The distinction this paper draws between a resolvable identifier and a local
asset has an analogue in data-management practice, where findability and
accessibility are separated from actual retrieval~\cite{wilkinson2016fair}. The
corpus used by the audited line of work is deposited with a permanent identifier
and is openly downloadable~\cite{locata,lollmann2018locata}, and it is
nonetheless absent here. What
Proposition~\ref{prop:offladder} adds is a reason not to fold that fact into the
ladder: an assertion that can hold when a lower rung fails, and fail when a
higher one holds, has no insertion position that preserves the demand
condition, so adding it would replace rather than refine the instrument.
Instruments that score availability on a single scale are not merely coarse;
they are measuring something that has no single scale.

\subsection{Summaries, outliers, and paired designs}

The second and third results of this paper join a much older discussion. That
the arithmetic mean is not resistant to a single extreme observation, and that
rank-based and median-based summaries are, is elementary~\cite{huber1964robust};
the paired signed-rank test we use throughout is one of the standard responses
to it~\cite{wilcoxon1945individual,dixon1946sign,dietterich1998approximate}, and
the intervals we report beside it are the ordinary bootstrap and binomial
constructions~\cite{efron1979bootstrap,clopper1934use}. What the literature has
less to say about is the reporting question, which is not which summary is
correct but what a reader can do when only one of them was printed.
Proposition~\ref{prop:disclosure} states the gap, and the record examined in
Section~\ref{sec:display} realises it: at a condition where every trial
completed in both arms, the mean and the median of the same forty pairs name
different estimators.

Localisation studies routinely report error against the Cram\'er--Rao bound as a
measure of how much of the available information an estimator recovers, and the
comparison is usually made through a root-mean-square error. That choice
inherits the sensitivity of the mean while appearing to be a property of the
estimator, and the efficiency analysis reported below shows what it costs on this record: a
factor of roughly eight between the two readings of the same trials against the
same bound.

Pre-registration and protocol-first reporting address a neighbouring failure, in
which the summary is chosen after the data are seen~\cite{nosek2018preregistration,
gelman2014garden,simmons2011false,amrhein2019retire}. We do not allege that here
and have no way to; the record we read declares its display rule in a field, and
declares the median primary. The problem we describe survives an honest
declaration, because the reader of a published table does not receive the trials
the declaration was applied to.

\subsection{Acoustic localisation, and what a reconstruction is asked to recompute}

The audited comparison sits in a line of work whose estimator is a time-delay
construction. The classical estimator of a delay from two sensors is a
generalised cross-correlation~\cite{knapp1976gcc}; later array methods and
challenge corpora built evaluation practice around that
family~\cite{allen1979image,lollmann2018locata,locata,
scheibler2018pyroomacoustics}. Signal-processing
venues were also among the first to treat a recomputable table as part of the
publication rather than as an optional
appendix~\cite{vandewalle2009repro}. The ladder below does not evaluate those
estimators. It asks a prior question: how far a reader on an ordinary machine
gets toward the table those papers would have the reader recompute.

\section{Methods}
%% Independent methods file. Not woven into main.tex prose.
%% Every quantity described here resolves to a hash-bound evidence entry.

\subsection{The ladder}

The instrument is a short ordered list of assertions, each of which can be
checked on a given machine without asking the original authors anything. In
increasing order of demand:

\begin{enumerate}
\item A citable published record for the work exists.
\item The reference source can be obtained at a fixed, named revision.
\item The runtime that the reference source requires is available on the machine.
\item The published table can be recomputed on the machine.
\end{enumerate}

Three properties make this more useful than a pass-or-fail verdict. It is
\emph{monotone}: a break at one rung makes every rung above it unreachable, so
the informative quantity is the index of the first break and not a count of
failures. It is \emph{machine-relative}: the same work can reach rung two on one
machine and rung four on another, and the ladder records which machine and when
rather than pretending to a universal property. And it is \emph{refutable}: each
rung names the observation that would settle it, so a reader with a different
machine can disagree with a specific line rather than with a conclusion.

Two things are stipulated not to satisfy a rung, because both are common and
both quietly convert a break into a false pass. A different interpreter that
executes a similar language is not the required runtime; source written for one
numerical environment may run, fail, or silently produce different results under
another, and treating the substitute as equivalent is an assumption, not a
reproduction. And a table produced by a re-implementation is not the published
table, however carefully the re-implementation was written.

Corpus availability is deliberately kept off the ladder and tracked as a separate
axis. It does not behave like the rungs: a corpus can be perfectly public,
openly licensed and permanently archived while still being absent from the
machine, and calling that state either reachable or broken loses the distinction
that matters. We record the public status and the local status as two separate
facts.

\subsection{The probe}

The project being audited already carried a recorded status for each of these
questions. That record is not what this study reports, because a recorded status
ages: a runtime can be installed, a corpus can be fetched, a pinned revision can
be moved. Instead a read-only probe re-checks each rung against the live machine
and writes what it observed with a timestamp, and the manuscript binds the probe
output.

The probe resolves the required interpreter and the common substitute on the
executable search path; confirms that the reference source directory exists,
reads its revision identifier, compares it against the recorded pin, and counts
how many of the recorded entry points are present; and checks a list of candidate
locations for the evaluation corpus, including the location named by the
project's own data-root environment variable. It installs nothing, downloads
nothing, and executes no vendor code. Its output records the first broken rung
directly, so the headline of the audit is a number the probe computed rather than
a reading of prose.

Reporting a probe of this kind is itself a design choice worth stating: an audit
whose evidence is a recollection cannot be re-run, and an audit that cannot be
re-run cannot be contradicted.

\subsection{The comparison that does remain reproducible}

The same project contains a fully local simulation study built on an open-source
room-acoustics package~\cite{scheibler2018pyroomacoustics}, and it supplies the
second half of this paper. Sources are localised from time-of-arrival
measurements at a microphone array in which each channel carries an unknown
constant timing offset. Two estimators are compared. The offset-agnostic
estimator solves an ordinary least-squares problem with all offsets fixed at
zero. The offset-aware estimator treats the offsets as unknown parameters with a
prior and returns a maximum a posteriori estimate over source position and
offsets jointly.

The independent variable is the dispersion of the true per-channel offsets across
\NConditions{} conditions, from a condition in which the offsets are exactly zero
up to a condition in which they are large relative to the measurement noise. Each
condition draws \NTrials{} trials, and both estimators are run on the same trial,
so the comparison is paired within trial. The reported quantity is the median
localisation error. That choice is inherited from the project's own display
rule rather than made here, and it matters: the error distribution has a heavy
upper tail, so the mean is pulled by outliers. Means are recorded in the
underlying artifacts and are not deleted, but they are not promoted to the
display quantity.

Each condition is tested with a two-sided Wilcoxon signed-rank
test~\cite{wilcoxon1945individual} from a standard scientific computing
library~\cite{virtanen2020scipy}, paired per trial. Alongside the test we report
the count of trials on which the offset-aware estimator was the better of the
two, because a count of \NTrials{} paired outcomes distinguishes a small
consistent advantage from a large but mixed one in a way a $p$ value does not.
The simulation was run once at a recorded seed and is not re-run here; this study
reads its output rather than regenerating it.

\subsection{Evidence discipline}

No number in this manuscript is typed by hand. Every artifact it rests on,
including the probe output, is recorded with its SHA-256 digest in a manifest
that accompanies the manuscript source. The figure and table code verifies each
digest against the bytes on disk before reading them and aborts on any mismatch,
and it emits both the figures and the numeric macros that the text prints. Even
the corpus size quoted above is extracted from the bound project record rather
than retyped, so the manuscript cannot disagree with the record it cites.


\section{Where the reconstruction stops}

\subsection{The observed ladder}

Fig.~\ref{fig:ladder} and Table~\ref{tab:ladder} give the probe's observation.
Rung one holds: the audited comparison has a citable record~\cite{hybridtdoa2024}.
Rung two holds and holds precisely: the reference source is present, its revision
identifier matches the pin recorded for it, and all \EntryPoints{} recorded entry
points are on disk. Nothing about the availability of this work is deficient, and
an artifact badge asserting availability would be accurate.

Rung \FirstBreak{} breaks. The source is written for a commercial numerical
computing environment that is not installed on this machine, and the free
interpreter most often proposed in its place is not installed either. We note the
second fact only for completeness, because as stated in the methods a substitute
interpreter would not have satisfied the rung: running source written for one
numerical environment under another is a claim about equivalence that would have
to be established, not assumed. Rung \NRungs{} is then unreachable, and the
audit's outcome is the integer \FirstBreak{} together with the observation
justifying it.

The project's own recorded state agrees, and it records \BlockedTiers{} distinct
blocked comparisons for three unrelated reasons: one blocked by the missing runtime,
one blocked by a dependency on both that runtime and the absent corpus, and one
blocked because the simulation retains only derived timing measurements and no
waveforms, so a whole family of signal-domain baselines cannot be run from the
stored data at all. That third blocker is instructive precisely because it has
nothing to do with licensing. Reachability fails for boring, various and specific
reasons, which is the argument for recording the specific reason rather than a
verdict.

\begin{table}[tbp]
\centering
\caption{The reachability ladder as observed on one machine on \ProbeDate{} by a
read-only probe. The ladder is monotone, so the first broken rung determines
every rung above it.}
\label{tab:ladder}
\small
% Generated by the figure script from hash-bound evidence. Do not hand-edit.
\begin{tabular}{clp{0.44\linewidth}c}
\toprule
Rung & Assertion & What was observed & Status \\
\midrule
1 & Citable record & A citable published record for the reference implementation exists. & reachable \\
2 & Source at a fixed revision & Source tree present at a fixed revision matching the recorded pin; all 19 recorded
entry points found. & reachable \\
3 & Required runtime present & The interpreter the source requires does not resolve on this machine. A commonly
suggested substitute is also absent, and would not count if present. & \textbf{broken} \\
4 & Published table recomputable & Not attempted. Unreachable through the rung below; no table is synthesised. & \textbf{broken} \\
\bottomrule
\end{tabular}

\end{table}

\begin{figure}[tbp]
\centering
\includegraphics[width=\linewidth]{fig1_reachability_ladder.pdf}
\caption{\CapLadder}
\label{fig:ladder}
\end{figure}

\subsection{Why the corpus is a separate axis}

The evaluation corpus used by this line of work~\cite{locata} is deposited in a
public archive with a permanent identifier and is roughly \CorpusGB{} gigabytes.
The probe
checked the locations where the project would keep it, including the one named by
its own data-root environment variable, and found nothing.

The temptation is to score this as reachable, on the grounds that anyone can
download it. Resisting that temptation is the point. A reproduction attempt that
requires tens of gigabytes of transfer, disk, and time is in a materially
different state from one that does not, and the difference is invisible to any
scheme that records only whether an asset is public. The two facts we record,
publicly archived and not locally present, are both true, and neither one implies
the other.

\section{What remains reproducible, and the condition it exposes}

\subsection{The expected result, and its exception}

The locally reproducible comparison behaves as the literature would predict over
most of its range. Table~\ref{tab:sweep} and Fig.~\ref{fig:sweep} show that as
the dispersion of the unknown per-channel offsets grows, the offset-agnostic
estimator degrades steadily, from a median error of \NaiveRangeLow{} m to
\NaiveRangeHigh{} m, while the offset-aware estimator stays nearly flat between
\AwareFloorLow{} m and \AwareFloorHigh{} m. At every one of the
\NonZeroConditions{} conditions with offsets present the paired signed-rank test
favours the offset-aware estimator.

The exception is the condition where the offsets are exactly zero. There the
offset-agnostic estimator is almost exact, with a median error of
\ZeroNaiveMedian{} m, while the offset-aware estimator returns
\ZeroAwareMedian{} m. The ratio is about \ZeroRatio{} to one and the paired test
gives $p = \ZeroP{}$ in the opposite direction to every other condition.

Fig.~\ref{fig:wins} makes the shape of that reversal clearer than the medians
can. Where offsets are present the offset-aware estimator wins a majority of
trials but never all of them, reaching \MaxOffsetWins{} of \NTrials{} at the
largest dispersion; the advantage is consistent, not absolute. Where offsets are
absent it wins \ZeroWins{} of \NTrials{}. A unanimous loss is a different and
more informative object than a significant one: it says the reversal is a
property of the estimator rather than a tail of its distribution.

\begin{table}[tbp]
\centering
\caption{Median localisation error by offset dispersion, over \NTrials{} paired
trials per condition at a fixed recorded seed. The paired difference is stated as
agnostic minus aware, so a positive value favours the offset-aware estimator. The
final columns give the count of trials on which the offset-aware estimator was
better and the two-sided paired signed-rank $p$ value.}
\label{tab:sweep}
\small
% Generated by the figure script from hash-bound evidence. Do not hand-edit.
\begin{tabular}{rrrrrc}
\toprule
Offset s.d. (ms) & Agnostic (m) & Aware (m) & Paired diff. (m) & Aware wins & Exact $p$ \\
\midrule
0 & 0.0015 & 0.5101 & -0.5086 & 0/40 & 1.82e-12 \\
5 & 0.8769 & 0.7219 & +0.0741 & 25/40 & 0.037 \\
10 & 1.1139 & 0.7345 & +0.4385 & 29/40 & 3.06e-05 \\
20 & 1.3186 & 0.7433 & +0.6301 & 36/40 & 1.81e-08 \\
30 & 1.3797 & 0.8352 & +0.6494 & 35/40 & 3.95e-08 \\
50 & 1.3900 & 0.6773 & +0.7272 & 29/40 & 8.15e-06 \\
75 & 1.6631 & 0.7070 & +0.8945 & 36/40 & 2.36e-08 \\
\bottomrule
\end{tabular}

\end{table}

\begin{figure}[tbp]
\centering
\includegraphics[width=0.88\linewidth]{fig2_offset_sweep.pdf}
\caption{\CapSweep}
\label{fig:sweep}
\end{figure}

\begin{figure}[tbp]
\centering
\includegraphics[width=0.88\linewidth]{fig3_paired_wins.pdf}
\caption{\CapWins}
\label{fig:wins}
\end{figure}

\subsection{Why the reversal happens, and why it is not a bug}

The mechanism is not mysterious and that is what makes it worth reporting. The
offset-aware estimator always estimates the per-channel offsets, because it
cannot know they are zero. Estimating parameters that carry no signal costs
variance, and that cost does not shrink as the true offsets shrink; it is a floor
the estimator pays for its own generality. The offset-agnostic estimator has no
such floor because it assumes the answer, and when the assumption is exactly
right the assumption is free. The near-flat aware curve and the near-exact
agnostic point at zero are two views of the same fact.

Read that way, the comparison is not evidence that one estimator is better. It is
evidence that the two estimators have different failure modes, and that the
experiment which distinguishes them is the one at the boundary of the modelled
effect. A study that swept only conditions with substantial offsets, which is the
natural design if the offsets are the phenomenon of interest, would have observed
a clean monotone advantage and reported it in good faith.

We think this generalises past acoustics into a small norm. When a method adds
parameters to absorb a nuisance, the condition worth reporting is the one where
the nuisance is absent, because that is where the price of the added generality
is visible and nowhere else. It is also the condition most likely to be dropped
from a sweep for looking uninformative.

\section{A comparison that runs, and still cannot be checked}
\label{sec:display}

\subsection{Three summaries, two answers}

Everything so far has been about whether a reader can get a comparison to run.
The third record in this project is one they can, and it is the more
uncomfortable case. It puts two locally written estimators for the same
calibration problem against each other at \CxConditions{} timing-noise levels,
\CxTrialsPerCondition{} paired trials each, and it stores both the cell summary
a paper would print and the \CxPlanned{} individual trials that summary came
from. That second half is what makes the following possible at all.

Fig.~\ref{fig:displayrule} and Table~\ref{tab:complement} read those trials
three ways. The median names the estimator given both measurement sets at every
condition. The paired test agrees at every condition, on
\CxPairedWinsLow{} to \CxPairedWinsHigh{} of the usable trials, at
$p \le \CxPairedPHigh{}$. The mean disagrees with both at \CxDisagreeing{} of
the \CxConditions{} conditions. Nothing separates the three except the
arithmetic applied at the end.

The cleanest of the three conditions is the one to read first, because it admits
no other explanation. At \CxCleanSigma{} microseconds of timing noise every
planned trial returned an error in both arms: nothing is missing, nothing was
dropped, no rule for handling failures is in play. The mean error is
\CxCleanMeanFull{} m against \CxCleanMeanReduced{} m, which names the
one-measurement estimator. The median is \CxCleanMedianFull{} m against
\CxCleanMedianReduced{} m, which names the other. The paired test says the
two-measurement estimator was better on \CxCleanWins{} of
\CxTrialsPerCondition{} trials at $p = \CxCleanP{}$. Two sentences, both true, both
computed from the same forty numbers, and they name different winners.

\subsection{One trial in forty}

Fig.~\ref{fig:pairedtrials} shows why, and it is the figure that makes the
argument rather than the tables. At every condition the cloud of paired trials
sits below the diagonal, which is the majority the paired test counts. Sitting
apart from the cloud, two orders of magnitude away, are a handful of trials on
which one estimator or the other went badly wrong. The mean is a statement about
those, and the median is a statement about the cloud.

Table~\ref{tab:influence} puts a size on it. The single largest trial of the
two-measurement estimator at the clean condition is \CxCleanLeverage{} times
that estimator's own median; removing it alone moves its mean from
\CxCleanMeanFull{} m to \CxCleanMeanWithoutMax{} m. Fig.~\ref{fig:influence}
turns that into a curve, and the curve is short. At every condition where the
mean disagrees with the median, setting aside \CxFlipDepth{} trial per arm --- out
of \CxTrialsPerCondition{} --- is enough for the mean to name the estimator the
median and the paired test had named all along, while the gap between the
medians barely moves.

We are not proposing that anyone trim. The point is the leverage: a conclusion
that one trial in \CxTrialsPerCondition{} can reverse is not a weaker version of
a conclusion that survives, it is a different kind of object, and the printed
row does not distinguish them.

\subsection{Two runs that produced nothing}

At two of the conditions one estimator diverged on one trial and returned no
error at all. The record writes those \CxAbsentValues{} values as the bare token
\texttt{NaN}, which the interchange format cannot express: a strict reader
rejects the whole file, and a permissive one accepts something number-shaped
that is not a number. Either way the failure stops being visible at the moment a
consumer opens the file, which is a small and entirely mundane mechanism by
which a summary loses the thing a reader most needs to know.

What the cell summary then reports is an average over what remained. The two
arms at those conditions are averaged over different numbers of trials, and the
printed row does not say so. By
Corollary~\ref{cor:completecase} that comparison is biased towards whichever arm
failed, since the dropped run is one that was going badly. Here the bias runs
against the conclusion we draw rather than for it: the arm that lost the trials
is the one the paired test finds worse, so the test reaches its answer despite
the drop. That is worth stating precisely because it is the direction in which
we would not have noticed.

\subsection{The same estimators against the same bound}

Fig.~\ref{fig:efficiency} makes the last point without reference to which
estimator wins. The Cram\'er--Rao bound for this problem is computed at each
trial from the array geometry and the timing noise; it is a property of the
problem, identical under either summary, and it is the natural yardstick for
asking how good either estimator is in absolute terms.

Read through the mean, the two estimators sit between \CxEffMeanLow{} and
\CxEffMeanHigh{} times the bound, and which one sits closer changes with the
noise level. Read through the median, they sit between \CxEffMedianLow{} and
\CxEffMedianHigh{} times it at every condition. The first reading says the
estimators are far from what the problem allows and behave erratically; the
second says both are essentially at the limit. The numerator is the only thing
that differs. A reader given the first number has no way to arrive at the
second, and no way to know that there is a second.

\subsection{What the disclosure would have cost}

The obvious objection is that per-trial records are burdensome, so it is worth
recording what this one weighs. The cell summaries in this record occupy
\CxBytesCells{} bytes. The two paired error columns --- every trial, both
estimators, full precision, and nothing else --- occupy \CxBytesMinimal{} bytes.
The disclosure that changes the answer is smaller than the summary that hides
it. The complete per-trial record, with every quantity each run reported and its
bound, is about \CxBytesTrials{} kilobytes, which is \CxCorpusOrders{} orders of
magnitude below the evaluation corpus this same project cannot fit on a laptop.
The three sizes are produced by one serialiser at full precision and are
comparable to each other rather than to any particular file format; the point is
the order of magnitude, and it does not depend on the format.

\begin{table}[tbp]
\centering
\caption{One record of \CxPlanned{} paired trials, read under \CxRules{}
summaries. Columns are timing-noise levels; the estimator named by each summary
is given in italics beneath the quantities that decide it. Trials returning an
error can be fewer than trials run, because a diverged run reports none.}
\label{tab:complement}
\small
% Generated by the figure script from hash-bound evidence. Do not hand-edit.
\begin{tabular}{lrrr}
\toprule
Timing-noise s.d. & 50\,$\mu$s & 100\,$\mu$s & 500\,$\mu$s \\
\midrule
Trials run per arm & 40 & 40 & 40 \\
Returned an error, both sets & 40 & 40 & 40 \\
Returned an error, one set & 40 & 39 & 39 \\
\midrule
Mean error (m), both sets & 0.1896 & 0.0762 & 0.8166 \\
Mean error (m), one set & 0.1217 & 1.0365 & 0.7258 \\
\quad \emph{named by the mean} & one set & both sets & one set \\
\midrule
Median error (m), both sets & 0.0257 & 0.0477 & 0.2359 \\
Median error (m), one set & 0.0312 & 0.0621 & 0.3396 \\
\quad \emph{named by the median} & both sets & both sets & both sets \\
\midrule
Paired trials usable & 40/40 & 39/40 & 39/40 \\
Won by both sets & 31/40 & 31/39 & 30/39 \\
Two-sided signed-rank $p$ & 2.17e-03 & 4.39e-04 & 0.018 \\
\quad \emph{named by the paired test} & both sets & both sets & both sets \\
\bottomrule
\end{tabular}

\end{table}

\begin{figure}[tbp]
\centering
\includegraphics[width=\linewidth]{fig4_display_rule.pdf}
\caption{\CapDisplayRule}
\label{fig:displayrule}
\end{figure}

\begin{figure}[tbp]
\centering
\includegraphics[width=\linewidth]{fig5_paired_trials.pdf}
\caption{\CapPairedTrials}
\label{fig:pairedtrials}
\end{figure}

\begin{table}[tbp]
\centering
\caption{The largest single trial of each estimator at each timing-noise level,
against that estimator's own median, and the mean recomputed without it. The
final column is the factor by which one trial inflates the printed mean.}
\label{tab:influence}
\small
% Generated by the figure script from hash-bound evidence. Do not hand-edit.
\begin{tabular}{llrrrrr}
\toprule
Noise s.d. & Given & Largest (m) & $\times$ own median & Mean (m) & Mean without it & Ratio \\
\midrule
50\,$\mu$s & both sets & 6.206 & 241 & 0.1896 & 0.0354 & 5.36 \\
 & one set & 2.177 & 69 & 0.1217 & 0.0690 & 1.76 \\
100\,$\mu$s & both sets & 0.973 & 20 & 0.0762 & 0.0532 & 1.43 \\
 & one set & 32.318 & 520 & 1.0365 & 0.2133 & 4.86 \\
500\,$\mu$s & both sets & 15.191 & 64 & 0.8166 & 0.4480 & 1.82 \\
 & one set & 7.482 & 22 & 0.7258 & 0.5480 & 1.32 \\
\bottomrule
\end{tabular}

\end{table}

\begin{figure}[tbp]
\centering
\includegraphics[width=\linewidth]{fig6_influence.pdf}
\caption{\CapInfluence}
\label{fig:influence}
\end{figure}

\begin{figure}[tbp]
\centering
\includegraphics[width=0.88\linewidth]{fig7_efficiency.pdf}
\caption{\CapEfficiency}
\label{fig:efficiency}
\end{figure}

\section{A condition that was never run}
\label{sec:breakeven}

\subsection{The reversal is an interval, not a point}

The offset study reports its reversal at a single condition, offsets exactly
zero, and a single condition is easy to set aside. No array is perfectly
synchronised, so a result that lives only at exactly zero describes a state no
deployment is in. The objection is fair and it is answerable, because the reason
the reversal looks like a point is visible in the design: between zero and the
lowest non-zero condition of \CoarseGapHigh{} milliseconds, that sweep has no
conditions at all.

We reran the comparison inside that interval. The estimators are the same code,
the trial count is the same \FineTrials{} paired trials, and the design is what
changed: \FineLevels{} dispersions instead of two, of which \FineAdded{} are new,
crossed with \FineRooms{} independently simulated arrays, for \FinePairs{} paired
trials in total. Panel A of Fig.~\ref{fig:breakeven} shows what is in
there. In every array the offset-agnostic estimator holds the lower median error
across most of the interval, at between \ReversalLevelsLow{} and
\ReversalLevelsHigh{} of the \FineNonzeroLevels{} non-zero dispersions, and up to
between \ReversalLow{} and \ReversalHigh{} milliseconds depending on the array.
The reversal is not a knife edge at zero. It is a region, and the coarse grid
skipped over it in one step.

\subsection{Where the sign changes, and where it has not}

Table~\ref{tab:breakeven} and Panel B of Fig.~\ref{fig:breakeven}
locate the change of sign. In \CrossedRooms{} of the \FineRooms{} arrays it falls
inside the refined grid, between \CrossingLow{} and \CrossingHigh{} milliseconds
--- a spread of \CrossingSpread{} milliseconds, which is tight enough to look
like a property of the estimator pair. In the other \CensoredRooms{} it has not
happened by the top of the grid, and we record those as censored rather than
assigning them the last condition we ran, for the reason
Definition~\ref{def:breakeven} makes explicit: a sweep that stops below a
crossing has not measured one.

That censoring is the finding, not a limitation of it. The top of our refined
grid is the lowest non-zero condition the original sweep ever ran. At that
condition, in \AnchorAgnosticRooms{} of \FineRooms{} arrays, the offset-agnostic
estimator still has the lower median error. A study that had drawn a different
room and started its grid where this one did would have read the opposite
headline off its own lowest condition, with the same code, the same trials per
condition and the same analysis.

The two summaries of the last section reappear here in a new place. Read through
the median of the trials and read through the count of paired wins, the same
trials put the same sign change between \RuleSepLow{} and \RuleSepHigh{}
milliseconds apart. In the last section the choice of summary decided which
estimator was named; here it decides where the boundary between them is. The
quantity has changed and the dependence has not.

\subsection{The break-even was derivable from the printed medians}

One number was fixed before the sweep ran. Interpolating the original sweep's
own printed medians between its two lowest conditions puts the break-even at
\PredictedCrossing{} milliseconds. In the array that sweep was run on, the
measured value is \MeasuredCrossing{} --- an error of \PredictionError{}
milliseconds, well inside the \PredictionTolerance{} millisecond tolerance we
agreed to in advance.

This is the sharpest form of the paper's argument that we can state. The missing
information was not expensive. It was not even missing: it is derivable from the
numbers the record already prints. A reader holding the coarse table could have
computed, to within a tenth of a millisecond, the condition whose absence made
the headline read the way it does. What they could not do is run it, because
Corollary~\ref{cor:onecondition} is not satisfiable by anything a reader has ---
not the summary, not the per-trial record, not the code. The trials at a
condition nobody ran do not exist.

We take some care not to over-read the agreement. The prediction succeeds for
the array the original sweep used, and \CensoredRooms{} of the other arrays have
no measured crossing to test it against; the rule extrapolated from one room is
not thereby established for rooms in general, and the honest summary is that the
coarse record locates its own missing condition and says nothing about anyone
else's.

\begin{table}[tbp]
\centering
\caption{Where the ordering between the two estimators changes sign, by array
and by summary rule, over a grid refined inside the interval the original design
skipped. The first pair of columns counts the non-zero dispersions at which the
offset-agnostic estimator holds the lower median error, and gives the highest of
them. A sign change recorded as a bound is one that had not occurred by the top
of the refined grid, which is the lowest non-zero condition the original sweep
ran; the final column names the estimator with the lower median error at exactly
that condition.}
\label{tab:breakeven}
\small
% Generated by the figure script from hash-bound evidence. Do not hand-edit.
\begin{tabular}{lccccc}
\toprule
& \multicolumn{2}{c}{Agnostic named} & \multicolumn{2}{c}{Sign change (ms)} & Named at \\
\cmidrule(lr){2-3} \cmidrule(lr){4-5}
Array & levels & up to (ms) & median & paired & \CoarseGapHigh{}\,ms \\
\midrule
as recorded & 7/11 & 4.0 & 4.03 & 2.44 & aware \\
seed 7 & 10/11 & 5.0 & $>$5.0 & $>$5.0 & agnostic \\
seed 999 & 9/11 & 4.0 & 4.21 & 2.88 & aware \\
seed 101 & 8/11 & 3.5 & 3.98 & 3.62 & aware \\
seed 2718 & 10/11 & 5.0 & $>$5.0 & $>$5.0 & agnostic \\
seed 31337 & 11/11 & 5.0 & $>$5.0 & $>$5.0 & agnostic \\
\bottomrule
\end{tabular}

\end{table}

\begin{figure}[tbp]
\centering
\includegraphics[width=\linewidth]{fig8_breakeven.pdf}
\caption{\CapBreakeven}
\label{fig:breakeven}
\end{figure}

The overlaid trajectories in Fig.~\ref{fig:breakeven} answer where each array
changes sign. They make a second question harder to see: at a fixed dispersion,
do the arrays agree about which estimator has the lower median error? Fig.~\ref{fig:robustness}
re-encodes exactly the same \FineRooms{} $\times$ \FineLevels{} cells as a
heatmap. The cell labels retain the signed difference and its metre unit, while
the colour makes the room dependence visible without asking a reader to trace
six curves. The pattern is descriptive of this finite grid only; it is not a
population estimate and the right-censored rooms remain right-censored in the
crossing table.

\begin{figure}[tbp]
\centering
\includegraphics[width=\linewidth]{fig9_robustness.pdf}
\caption{\CapRobustness}
\label{fig:robustness}
\end{figure}

The finite-grid result has a second, deliberately different sensitivity check.
Fig.~\ref{fig:geometry-sensitivity} treats each of the twenty independently
drawn room/source/microphone geometries as one observation and resamples those
geometries, rather than pooling their trials. The signed advantage of the
offset-aware estimator grows across the tested dispersion range; the labels
give the number of geometries in which it wins. This is stronger evidence than
a trial-level interval, but it remains a local simulator result and not a
claim about an acoustic population.

\begin{figure}[tbp]
\centering
\includegraphics[width=\linewidth]{fig10_geometry_sensitivity.pdf}
\caption{\CapGeometrySensitivity}
\label{fig:geometry-sensitivity}
\end{figure}

\section{Discussion}

A reachability audit is cheap. The probe used here runs in well under a second,
installs nothing and touches no network. What it buys is a statement with a
falsifiable shape: not that the work is irreproducible, but that on a stated
machine on a stated date the first break was at a stated rung for a stated
reason. Anyone with a different machine can produce a different integer, and the
disagreement is then about an observation rather than about a judgement.

The audited work is not deficient in any way this instrument can detect other
than its runtime requirement, and that requirement is a property of the
scientific software ecosystem rather than of the authors. Writing numerical work
for a widely licensed commercial environment was, and remains, an ordinary
choice. The point of recording the break is not to attribute fault; it is that a
reader deciding whether to attempt a reproduction currently has no cheap way to
learn that a licence stands between them and rung four, and that the badge on the
artifact will not tell them.

The four parts of this paper are connected by more than a shared project, and
the order they arrive in is the argument. The reachability audit says which
comparisons a reader can check. The offset study says what a reader finds when
they check one: a headline that holds under the conditions its authors would
naturally have swept, and reverses under the one they would naturally have
omitted. The third comparison says that checking is not finished when the code
runs, because the summary a reader is handed does not determine the conclusion
the trials support. The refined sweep says that disclosing the trials is not the
end either, because a comparison also has conditions, and the ones a design
skipped leave no trace in the ones it kept.

Those last two are worth separating, because they ask for different things. A
missing per-trial record is recoverable: the runs happened, the bytes exist, and
releasing them costs less than the table they would accompany. A missing
condition is not recoverable at any price a reader can pay. It has no trials, no
summary and no residue in the record, and the only repair is to run it ---
which, by Corollary~\ref{cor:onecondition}, can be one condition when it is
placed at an unobserved point where the sign is negative. The uncomfortable part
of Section~\ref{sec:breakeven} is
that the record examined here did know: the position was computable from its own
printed medians, to within a tenth of a millisecond, by anyone who thought to
ask.

The point about summaries is the one that changes what we would build next. A ladder whose
top rung is ``the published table recomputes'' treats reproduction as the end of
the road, and Proposition~\ref{prop:disclosure} says it is not: the rung
certifies a number, and a claim is a number plus a rule for reading it. The
missing axis is disclosure of the per-trial record, and the reason to insist on
it is not that summaries are dishonest. The record examined here declares its
display rule in a field, names the median primary, and annotates the mean as
outlier-sensitive. Its authors did the careful thing. A reader holding only the
printed table still cannot check the conclusion, because the declaration travels
without the trials it was applied to.

Disclosure is also the cheapest of the three things this paper asks for. A
runtime licence costs money, a corpus costs tens of gigabytes, and the paired
error columns for \CxPlanned{} trials cost \CxBytesMinimal{} bytes, which is
less than the summary table that would sit beside them. We do not think the
obstacle to per-trial reporting is cost, and the numbers here make it hard to
argue that it is.

\section{What this study does not claim}

No claim of novelty in acoustic localisation is made, and no comparison here is
a finding about a published method. The estimators in the second and third
studies are local implementations written inside this project; they are named by
what they are given rather than by whose method they resemble, and neither is
anyone's released code. No substitute runtime is treated as equivalent to the
required one, no published table is reconstructed, approximated or synthesised,
and rung \NRungs{} is reported as unreachable and left unreached.

No permanent property of the audited work is asserted. The audit is
machine-relative and dated: installing the required runtime, or fetching the
corpus, changes the result, and the correct reading of a first break at rung
\FirstBreak{} is a statement about a machine on a day. No misconduct or
negligence is alleged, and nothing here assesses the quality of the audited
work, which this study did not evaluate.

The refined sweep is a new run, made for this manuscript, and it is not part of
the record being audited. It changes nothing in that record and reads nothing
back into it: the estimator code is imported unmodified, and the sweep is stored
as a separate artifact beside the existing ones. It is not part of any
pre-registered plan for this study either, and is reported as an extension
rather than as a confirmation of anything stated in advance. What was fixed in
advance is narrower and is stated where it is used: the prediction the refined
sweep is read against, and the tolerance for calling it successful, were both
computed and written down before the sweep ran.

Nothing in Section~\ref{sec:breakeven} is a claim about acoustic localisation or
about how anyone should have designed a sweep. A grid placed where the modelled
effect is expected to matter is a reasonable grid, and the point is not that its
authors erred but that the sign it reports is not identified by the estimators
alone. The break-even values are properties of these two implementations in
simulated rooms at a fixed prior width and noise floor; they are not a
recommendation for any array, and no claim is made that the prediction rule
which located one array's crossing would locate another's, since
\CensoredRooms{} of the arrays here have no measured crossing to test it
against.

The mean is not promoted anywhere in this paper, and neither is it discarded.
Section~\ref{sec:display} reports what the mean says in order to show that it
says something different from the trials it summarises, which is an argument for
disclosing the trials rather than for choosing between summaries. No result
reported here is computed from a trimmed sample; the trimming in
Fig.~\ref{fig:influence} measures leverage and is not applied to any quantity
stated as a finding.

\section{Conclusion}

A published comparison was audited against a \NRungs{}-rung reachability ladder
by a read-only probe. Its source is obtainable at exactly the revision recorded
for it and all \EntryPoints{} entry points are present, so availability is not
the problem; the first break is at rung \FirstBreak{}, the absence of the
commercial numerical environment the source requires, and the published table is
consequently out of reach. Corpus availability was tracked separately and shows
why it must be: the evaluation corpus is public, permanently archived and about
\CorpusGB{} gigabytes, and is absent locally, which is neither a satisfied nor a
broken requirement.

The reproducible remainder of the same project produced two further results.
Over \NConditions{} conditions at \NTrials{} paired trials each, an offset-aware
estimator beat an offset-agnostic one wherever the offsets it models existed,
and lost to it by a factor of about \ZeroRatio{} on every single trial of the
one condition where they did not. And in a third comparison of \CxPlanned{}
paired trials, the mean and the median of the same trials named different
estimators at \CxDisagreeing{} of \CxConditions{} conditions, including one at
which nothing was missing; a single trial in \CxTrialsPerCondition{} carried the
difference, and the same trials read against the Cram\'er--Rao bound put the
estimators at \CxEffMeanLow{} to \CxEffMeanHigh{} times it or at
\CxEffMedianLow{} to \CxEffMedianHigh{} times it depending only on which summary
is used.

A fourth result asked where that reversal ends. Rerunning the same estimators on
a grid refined inside the interval the original design stepped over, crossed
with \FineRooms{} independently simulated arrays and \FinePairs{} paired trials,
puts the offset-agnostic estimator ahead out to between \ReversalLow{} and
\ReversalHigh{} milliseconds of dispersion depending on the array, and still
ahead at the original sweep's lowest non-zero condition in
\AnchorAgnosticRooms{} of \FineRooms{} arrays. The sign a sweep reports is
therefore fixed by its grid as much as by its estimators, and the position of
the missing condition was computable from the original record's own printed
medians to within \PredictionError{} milliseconds.

The four results point the same way: the information a reader needs is the
specific place where something stops holding, and a reachability verdict, a
method headline, a summary statistic and a grid each hide exactly that. Three of
the four are cheap to fix. Record the rung where a reconstruction stops rather
than a badge; publish the per-trial record beside the table, which here would
have cost \CxBytesMinimal{} bytes, less than the table itself; and add the one
condition on the far side of the break-even, which costs one row and which the
table itself will tell you where to put.

\appendix

\section{Source artifacts and their checksums}
\label{app:evidence}

Table~\ref{tab:evidence} lists all \NEvidence{} artifacts behind this
manuscript, totalling \EvidenceBytes{} bytes, in the order the analysis consumes
them. Every number the text prints is computed from these bytes, and the code
recomputes each checksum before reading, so a stale artifact stops the build
rather than reaching the page.

This paper needs the discipline more than most, because its central claim is
about what could \emph{not} be reconstructed on one machine on one day. A claim
of that shape is unfalsifiable unless the state it describes is recorded well
enough for someone else to disagree with it. The read-only re-check of every
rung is therefore recorded here alongside the measurements, and it is listed as
what it is: an artifact produced for this manuscript, not an upstream record.
The reader is entitled to know which rows are the project's own testimony and
which are prior artifacts it merely read, and the descriptions say so.

Two limits keep the table honest. A checksum fixes bytes to a claim; it does not
make the bytes true, and none of this certifies that an artifact was correct
when written. And each entry's internal identifier and path are omitted
deliberately --- they describe a private tree, they would tell a reader nothing,
and keeping them out of the PDF is a property the build verifies rather than a
courtesy. That omission matters here for a further reason: this manuscript
audits a named third-party artifact, and the fewer incidental internal strings
it carries about that artifact, the smaller the chance of misattributing
something to it.

\begin{table}[htbp]
\centering
\caption{Every artifact this manuscript rests on, with its size and the first 16
hexadecimal characters of its SHA-256 checksum. Full checksums, and the paths they
apply to, are in the manifest accompanying the manuscript source; the prefix is
shown here only because a 64-character checksum does not set in a table column.
The figure and table code recomputes each checksum from the bytes on disk before
reading them and stops the build on any mismatch.}
\label{tab:evidence}
\footnotesize
% Generated by the figure script from hash-bound evidence. Do not hand-edit.
\begingroup
\setlength{\tabcolsep}{5pt}
\begin{tabular}{@{}p{0.56\linewidth}rl@{}}
\toprule
What it records & Bytes & SHA-256 (first 16) \\
\midrule
read-only re-check of every ladder rung in the local environment, timestamped; produced for this manuscript, not an upstream artifact &   1,839 & \texttt{5e2ad88b5c80696b} \\
7 offset conditions x 40 paired trials, per-condition median and mean error and paired signed-rank &   4,941 & \texttt{bea4b6733963dadd} \\
extracted paired test per condition, with the analysis procedure, seed and timestamp of the sweep &   3,557 & \texttt{5c5f2892fb74b72f} \\
what the project itself records as present, blocked, and why &   1,934 & \texttt{a08fc117afabc171} \\
reference-implementation entry points, revision pin, and the stated execution blocker &   1,253 & \texttt{42498d6f993ddc1d} \\
the published-table tier, recorded as blocked &     377 & \texttt{81fd2de7dab54310} \\
the full-replication tier, recorded as blocked on both runtime and corpus &     304 & \texttt{ee61d33aa06acfa9} \\
a fourth tier blocked for an unrelated reason, kept to show the ladder is not one failure &     585 & \texttt{a1c9a9f0fcb6a96e} \\
summary of the reproducible baseline arm at one condition &     383 & \texttt{2f7faa2e35672d5d} \\
summary of the reproducible treatment arm at the same condition &     408 & \texttt{17910cc8cc0834ae} \\
independent re-implementation run, explicitly flagged as not the published table & 148,676 & \texttt{d240f0308b9526fb} \\
a provenance record of display-metric rules, blockers and the archive disposition &   3,943 & \texttt{9a79afd5808186ad} \\
refined dispersion grid inside the interval the coarse sweep skipped, crossed with 6 geometries, 40 paired trials per cell, every trial stored; run for this manuscript & 186,240 & \texttt{2929cb29383960fd} \\
per-cell summaries, the located sign change under two summary rules, and the check of the prediction the coarse record fixes in advance &  58,959 & \texttt{c4d2e7891617aa1a} \\
geometry-level bootstrap sensitivity audit over the existing 20 independent room/source/microphone geometries; resampling unit is geometry, not trial &   2,100 & \texttt{9bd14f721b452cbc} \\
microphone-count x offset sensitivity sweep; descriptive-only robustness readout &  28,440 & \texttt{725b7db4c48c3c0e} \\
reverberation-severity x offset sensitivity sweep; descriptive-only robustness readout &  17,671 & \texttt{352d1567c8b6937c} \\
joint-posterior calibration case study with SBC coverage and divergence diagnostics &   9,786 & \texttt{3fff057fb9115c33} \\
\bottomrule
\end{tabular}
\endgroup

\end{table}

\section*{Data and code availability}

The probe, the derived tables, the recorded project state and the figure code
that produces every number in this manuscript accompany the manuscript source.
Each artifact is recorded with its SHA-256 checksum in a manifest, and the figure
code verifies those checksums against the bytes on disk before reading them, so the
build fails rather than silently reporting stale numbers. The reference source
examined in the audit is third-party software, was not modified and is not
redistributed here.
%% Generated by scripts/release_targets.py from the release ledger.
%% Do not edit: an identifier typed here would not be checked against
%% anything, which is exactly the failure the ledger exists to prevent.
%% Emitted 2026-08-28T12:16:18Z
No public repository has been created for that material and no persistent identifier has been issued for it, so neither is quoted here; both will be recorded when they exist and not before.


\section*{Ethics statement}

This study is entirely computational and involves no human participants, no
animal subjects and no personally identifying data. It reports on the
reachability of third-party research software on one machine. It makes no
assessment of the correctness or quality of that software and alleges no
wrongdoing.

\section*{Competing interests}

The authors declare no competing interests.

\section*{Funding}

No external funding was received for this study.

\section*{Author contributions}

\ifdefined\ANONYMOUS
The sole author performed conceptualization, methodology, software, validation,
formal analysis, investigation, visualization, writing---original draft, and
writing---review and editing.
\else
Zeyu Fu: conceptualization, methodology, software, validation, formal analysis,
investigation, visualization, writing---original draft, and writing---review
and editing.
\fi

\section*{Declaration of generative AI use}

Generative AI tools assisted with code preparation and language editing. Every
number printed above is emitted programmatically from artifacts tied to the manuscript by cryptographic
checksum rather than transcribed into prose. Structural claims are
enforced by build-time checks that stop the build rather than by unverified
sentences. Responsibility for the content rests with the authors.
